\section*{Open questions (5, with concrete next steps)}

\begin{enumerate}

\item \textbf{Is PprS:pprM binding direct or scaffolded?}
\emph{Basis.} The mechanistic claim rests on MAPS pull-down + EMSA + Northern. MAPS is not base-pair-resolved; EMSA uses purified RNA and cannot exclude in-cell bridging factors. No CLIP-seq / CLASH / PARE-seq footprint is deposited. My redo confirms MAPS enrichment quantitatively ($\rho=0.89$, \emph{pprM} in top decile) but cannot address directness.
\emph{Next steps.} Run CLASH or in vivo psoralen-crosslink RNA--RNA ligation (COMRADES / PARIS) on \emph{D.\ radiodurans} WT vs PprSKD at 10 kGy IR to recover the exact PprS:\emph{pprM} chimera and pairing footprint. Compare to the IntaRNA in silico prediction. Direct footprint at the predicted \emph{pprM} 5$'$UTR site confirms mechanism; no chimera despite high MAPS enrichment means the interaction is indirect and the model needs revision.

\item \textbf{Does the PprS/pprM axis transfer across the Deinococcus genus?}
\emph{Basis.} The paper studies only \emph{D.\ radiodurans} R1. The genus spans $\sim$10--100$\times$ IR-tolerance; a causal PprS/\emph{pprM} axis should give a dose--response with radioresistance level. Comparative context is absent from the paper and my reproduction.
\emph{Next steps.} BLASTn PprS (Dsr2) and \emph{pprM} against \emph{D.\ geothermalis} DSM 11300, \emph{D.\ deserti} VCD115, \emph{D.\ proteolyticus} MRP; where homologs exist, construct markerless PprSKD; assay 10 kGy IR survival vs WT for each pair; compare KD-induced sensitization slope vs WT IR-D10. Monotone dose--response = genus-wide mechanism; R1-only effect = strain-specific rewiring.

\item \textbf{Is PprS/pprM coupled to the Mn-antioxidant defense arm, or is it a repair-only track?}
\emph{Basis.} \emph{D.\ radiodurans} radioresistance is jointly attributed to efficient DNA repair \emph{and} high intracellular Mn:Fe ratio protecting proteins from oxidative damage. The paper treats PprS/\emph{pprM} as pure repair regulation and does not cross-reference the Mn-peptide pool. My IR-effect gene lists (S6/S7) contain repair genes but I did not check enrichment of Mn-transport (MntH, MntABC), Mn-SOD (SodA), or Mn-peptide biosynthesis loci.
\emph{Next steps.} Measure intracellular Mn:Fe ratio and Mn-peptide pool (small-peptide LC-MS) in WT vs PprSKD at 0 vs 10 kGy; reduced Mn-peptide loading in PprSKD implies coupling. Also re-mine my per-gene TSVs (all 3132 features present) against a curated Mn-metabolism gene set; no new experiments needed for the in-silico half.

\item \textbf{Is the failed S4 contrast recoverable with R-DESeq2, or is the paper's headline padj a specific-parameterization artifact?}
\emph{Basis.} PyDeseq2 0.5.4 reproduces S6/S7 perfectly on the same public counts but gives Spearman = 0.024 on S4 with \emph{pprM} padj drifting from $1.3\times10^{-9}$ to $0.52$. Six S4 genes are tied at padj $= 6.28\times10^{-14}$ (a BH-floor pattern suggesting unusual size-factor behavior). WT-0 CV = 0.51 in library size is a textbook DESeq2 trouble spot at n = 3.
\emph{Next steps.} Install R + Bioconductor + DESeq2 on uicgpu; re-run S4 from the same 12 htseq files under (i) DESeq2 defaults, (ii) design \texttt{$\sim$strain+dose+strain:dose} tested via the interaction term, (iii) $\pm$ Cook's outlier refit, (iv) $\pm$ \texttt{lfcShrink} apeglm. R-DESeq2 recovering the paper's padj under some parameterization = PyDeseq2 numerical drift (bug report). R-DESeq2 also giving padj $\sim 0.5$ on default = paper's S4 list is a specific-parameterization artifact and the \emph{pprM} headline needs a caveat.

\item \textbf{What is the full sRNA-target regulatory network in \emph{D.\ radiodurans} (all 31 Dsr* sRNAs)?}
\emph{Basis.} The paper isolates one sRNA (PprS/Dsr2) and one target (\emph{pprM}). The htseq count tables show 31 annotated Dsr* sRNAs with quantifiable expression in all 12 RNA-seq samples. Neither the paper nor my reproduction attempts a genome-scale sRNA-target network reconstruction, but the raw data support it.
\emph{Next steps.} Run MAPS pull-down on the other 30 Dsr* sRNAs (or the top 5 most IR-responsive by S6/S7 change) to enumerate sRNA-target edges genome-wide; combine with in silico IntaRNA / TargetRNA3 for prioritization; cluster the resulting bipartite graph and check whether modules align with GO categories (DNA repair, Mn transport, translation, cell envelope). Feasible as a $\sim$2-month campaign at a UT Austin-scale lab. Deliverable: first sRNA-target regulatory atlas for \emph{D.\ radiodurans}, generalizing the \emph{pprM} finding to a network-scale claim.

\end{enumerate}
